//Falar sobre a estrutura do projeto, como ele é organizado, quais são os principais componentes, etc. a cena de este estar dividido em containers, o nginx as networks etc etc
\chapter{Project Structure} \label{ch:project_structure}

\section{System Architecture and Infrastructure}

The Hubly project is built upon a containerized infrastructure designed to ensure consistency across development and production environments, while maintaining a robust security posture. By leveraging \textbf{Docker}, we encapsulate each component—Frontend, Backend, and Database—into autonomous, isolated services, effectively eliminating environmental discrepancies.

\subsection{Containerization Strategy}
The infrastructure relies on \textit{Docker Compose} to orchestrate the lifecycle and communication between four core services:
\begin{itemize}
    \item \textbf{hubly-web}: A Node.js container serving the frontend application, optimized for hot-reloading during the development phase.
    \item \textbf{hubly-api}: A .NET 8 container hosting the backend logic, configured to interact securely with the persistence layer.
    \item \textbf{hubly-db}: A PostgreSQL 16 container, serving as the relational database.
    \item \textbf{hubly-nginx}: A lightweight Nginx container acting as a \textit{Reverse Proxy}.
\end{itemize}

\subsection{Network Segmentation and Security}
To enhance system security, the infrastructure implements a \textit{Strategic Network Segmentation} approach, creating two distinct virtual networks:

\begin{itemize}
    \item \textbf{app-net}: A public-facing network that facilitates communication between the Nginx proxy, the Frontend, and the Backend.
    \item \textbf{db-net}: A private, isolated network dedicated exclusively to Backend-Database communication.
\end{itemize}

By configuring the database service to reside solely within the \texttt{db-net} segment, we enforce a strict \textit{Enhanced Security Layer}. This prevents any direct external access to the database, ensuring that all data persistence requests must pass through the Backend’s authentication and validation filters.

\subsection{Reverse Proxy Configuration}
The Nginx container serves as the primary entry point for all incoming requests. Its configuration acts as a traffic manager, directing requests to either the API (via the \texttt{/api} route) or the Frontend application. This setup allows for centralized management of headers, SSL/TLS termination, and WebSocket upgrades, which are essential for supporting the real-time features provided by SignalR.

\subsection{Advantages of Containerized Distribution}

The decision to distribute our infrastructure into multiple specialized containers offers several critical advantages that enhance the scalability, maintainability, and security of the Hubly project:

\begin{itemize}
\item \textbf{Environmental Consistency}: By encapsulating each service within its own container, we ensure that the runtime environment is identical across all development, testing, and production stages. This eliminates the "it works on my machine" phenomenon, guaranteeing that dependencies and configurations remain consistent regardless of the underlying host.
\item \textbf{Decoupled Scalability}: Since each component (Frontend, Backend, Database) operates as an independent unit, we can scale specific parts of the application according to demand. For instance, the Backend can be scaled to handle higher computational loads without requiring changes to the Frontend or the database infrastructure.
\item \textbf{Fault Isolation}: The distributed nature of our architecture acts as a natural barrier against cascading failures. If a service experiences a runtime error or instability, the remaining containers continue to operate independently, ensuring that the system remains partially available and resilient to localized issues.
\end{itemize}

\subsection{Strategic Network Segmentation and Security}

Security is a foundational pillar of our architectural design, achieved primarily through rigorous \textbf{Network Isolation}. By leveraging Docker's bridge network driver, we have created a "zero-trust" environment between our internal services:

\begin{itemize}
\item \textbf{Hardened Perimeter}: The Frontend and Nginx are placed within the \texttt{app-net} network, acting as the public-facing layer. The database, however, resides exclusively within the private \texttt{db-net} network.
\item \textbf{Access Control}: By ensuring that the \texttt{hubly-db} container is not reachable via the \texttt{app-net} network, we physically and logically isolate the persistence layer. Any access to the data must originate from the \texttt{hubly-api} container, which is the only service permitted to span both networks.
\item \textbf{Attack Surface Reduction}: This isolation significantly reduces the application's attack surface. Even if an external actor were to compromise the Frontend or the Nginx proxy, they would lack the network routing necessary to establish a direct connection to the database. Consequently, all data access is strictly gated by the Backend’s authentication and validation \textit{pipeline}, preventing unauthorized or direct database exploitation.
\end{itemize}